\documentclass[profile=standard]{ipara-handbook}
\title{B09｜协方差矩阵与二次型：随机波动的方向结构}
\author{}
\date{}
\begin{document}
\maketitle

\begin{quote}
本册回答：\textbf{协方差矩阵为什么一定半正定？相关系数为什么不能任意超过 $1$？二维正态的“椭圆方向”为什么会由协方差矩阵的特征向量决定？}
\end{quote}

\section{本册定位：把概率的二阶摘要翻译成线代的方向二次型}
概率 Topic05 Own 方差、协方差；概率 Topic04 Own 二维正态与随机向量线性变换；线代 Topic05/06 Own 实对称谱结构、二次型与正半定。B09 只 Own：
\[
\boxed{
\Sigma
\xrightarrow{a^T(\cdot)}
\operatorname{Var}(a^TX)=a^T\Sigma a
\xrightarrow{\ge0}
\Sigma\succeq0
}
\]
以及线性变换下
\[
\boxed{\operatorname{Cov}(AX)=A\Sigma A^T.}
\]

\section{协方差矩阵不是一堆成对公式}
对随机向量
\[
X=(X_1,\dots,X_n)^T,
\qquad \mu=E[X],
\]
定义
\[
\Sigma=E[(X-\mu)(X-\mu)^T].
\]
矩阵元素是
\[
\Sigma_{ij}=\operatorname{Cov}(X_i,X_j).
\]
它首先是一张实对称矩阵，但更重要的是：它把所有方向上线性组合的波动统一编码起来。

\section{母接口：任意方向方差就是一个二次型}
取任意确定向量 $a$，标量随机变量
\[
Y=a^TX
\]
的方差为
\[
\begin{aligned}
\operatorname{Var}(Y)
&=E\bigl[(a^T(X-\mu))^2\bigr]\\
&=E\bigl[a^T(X-\mu)(X-\mu)^Ta\bigr]\\
&=a^T\Sigma a.
\end{aligned}
\]
因为任何方差都非负，故
\[
\boxed{a^T\Sigma a\ge0,\qquad \forall a.}
\]
于是
\[
\boxed{\Sigma\succeq0.}
\]

所以“协方差矩阵半正定”不是额外记忆性质，而只是“任意线性组合的方差不能为负”翻译到线代之后的必然结果。

\section{二维相关系数界为什么来自正半定}
二维协方差矩阵为
\[
\Sigma=
\begin{pmatrix}
\sigma_X^2&\rho\sigma_X\sigma_Y\\
\rho\sigma_X\sigma_Y&\sigma_Y^2
\end{pmatrix}.
\]
半正定要求
\[
\det\Sigma
=\sigma_X^2\sigma_Y^2(1-\rho^2)\ge0.
\]
若两方差均正，则
\[
\boxed{|\rho|\le1.}
\]

因此相关系数界也可以从二次型合法性重新生成，而不必把 Cauchy--Schwarz 与矩阵判据当成两套无关结论。

\section{线性变换：随机方向怎样被重新混合}
令
\[
Y=AX+b.
\]
均值平移为
\[
E[Y]=A\mu+b,
\]
而中心化后
\[
Y-EY=A(X-\mu).
\]
所以
\[
\boxed{
\operatorname{Cov}(Y)
=A\Sigma A^T.
}
\]

这与二次型变量代换非常相似：一个方向 $c$ 在新变量中的方差为
\[
\operatorname{Var}(c^TY)
=c^TA\Sigma A^Tc
=(A^Tc)^T\Sigma(A^Tc).
\]
也就是说，新坐标中的某个方向，实际对应原随机空间中的方向 $A^Tc$。

\section{正交主轴：特征向量为什么是主波动方向}
由于 $\Sigma$ 实对称且半正定，存在正交矩阵 $Q$ 使
\[
Q^T\Sigma Q
=\operatorname{diag}(\lambda_1,\dots,\lambda_n),
\qquad \lambda_i\ge0.
\]
令
\[
Z=Q^T(X-\mu).
\]
则
\[
\operatorname{Cov}(Z)
=Q^T\Sigma Q
=\operatorname{diag}(\lambda_i).
\]
所以在这些正交方向上，交叉协方差被消掉，每个坐标的方差就是对应特征值。

这给出一个重要几何读法：
\[
\boxed{\text{特征向量 = 二阶波动的主方向，特征值 = 该方向上的方差强度}.}
\]

注意这里只能一般推出“不相关”，不能推出独立；联合正态时才有额外加强。

\section{退化：零特征值不是“方差为零一点”，而是一条确定关系}
若存在非零 $a$ 使
\[
a^T\Sigma a=0,
\]
则
\[
\operatorname{Var}(a^TX)=0.
\]
因此
\[
a^TX=E[a^TX]
\]
几乎处处，即某个非平凡线性组合实际上没有随机波动。

二维中 $|\rho|=1$ 就对应协方差矩阵退化；两个中心化变量存在严格线性关系。这也解释非退化二维正态为什么要求
\[
|\rho|<1.
\]

\section{二维正态：概率密度中的二次型从哪里来}
非退化二维正态密度的指数核心可以写成
\[
(x-\mu)^T\Sigma^{-1}(x-\mu).
\]
这是一个正定二次型。沿协方差矩阵的正交特征方向换坐标后，若
\[
\Sigma=Q\Lambda Q^T,
\]
则
\[
(x-\mu)^T\Sigma^{-1}(x-\mu)
=z^T\Lambda^{-1}z
=\sum_i\frac{z_i^2}{\lambda_i}.
\]
所以等密度曲线的主轴方向由 $\Sigma$ 的特征向量给出，而轴向尺度由特征值决定。

这正是线代二次型与概率二维正态的稳定接口，而不是“椭圆看起来像二次型”的表面类比。

\section{B09 串起哪些章节}
\begin{tabular}{p{0.23\linewidth}p{0.30\linewidth}p{0.39\linewidth}}
\hline
Owner / 下游 & 提供对象 & B09 的翻译责任\\
\hline
概率 Topic05 & 方差、协方差、相关系数 & 把所有线性组合方差统一成二次型\\
概率 Topic04 & 二维正态、随机向量线性变换 & 解释 $A\Sigma A^T$、退化与等密度主轴\\
线代 Topic06 & 正半定二次型 & 把“方差非负”翻译为矩阵半正定\\
线代 Topic05 & 实对称谱定理 & 把协方差的方向耦合解成正交主波动方向\\
\hline
\end{tabular}

\section{反桥梁}
\begin{tabular}{p{0.20\linewidth}p{0.04\linewidth}p{0.20\linewidth}p{0.48\linewidth}}
\hline
A & & B & 真正区别\\
\hline
零协方差 & $\neq$ & 独立 & 前者只消除二阶线性交叠；一般联合结构仍可能依赖\\
协方差矩阵半正定 & $\neq$ & 正定 & 半正定允许零方差线性组合和低维退化\\
特征方向不相关 & $\neq$ & 特征方向独立 & 只有联合正态等额外结构下才可升级\\
协方差矩阵 & $\neq$ & 完整联合分布 & 它只是二阶摘要，很多不同分布共享同一协方差\\
\hline
\end{tabular}

\section{调用信号与一页压缩}
看到“$aX+bY$ 的方差、协方差矩阵、相关系数合法范围、二维正态线性变换、退化、等密度椭圆”时，先写
\[
\boxed{\operatorname{Var}(a^TX)=a^T\Sigma a.}
\]
然后问：
\begin{enumerate}
\item 当前方向 $a$ 是什么？
\item $\Sigma$ 是否半正定/退化？
\item 是否需要正交对角化读取主波动方向？
\item 当前结论只到“不相关”，还是题目额外给了联合正态资格？
\end{enumerate}

\end{document}
